\section{Further research questions}
\begin{enumerate}[leftmargin=1.5em]
\item Is there an effective algorithm, for all rank-two integral bouquet
voltages, that computes both the eventual polynomial and its least starting
level? The examples suggest reducing the symbol in characteristic $p$,
then handling finitely many exceptional directions over a small cyclotomic
field. An algorithm needs a termination proof and an effective bound on
the precision in those directions.
\item For $p>2m$ in dimension two, Theorem~\ref{thm:moments} reduces the
extremal problem to the orders of binary moment forms. Determine the largest
$A$ when the quadratic form vanishes identically. In higher dimension,
can one express the distribution of $A$ for random step lists through
point counts of the successive moment varieties?
\item For fixed $q$, determine the rate of decay of
$\Delta(\mathbb F_q^b)$. The density Hales--Jewett reduction gives
vanishing and the Behrend construction gives a subexponential-in-$\log b$
lower bound. Can additional polynomial ratios beyond those with entries
$1,x$ substantially sharpen the upper bound?
\item Classify growing products of finite local rings when their radicals
are nonzero. The reduced-ring proof uses coordinate degrees in integral
domains for the lower bound. Nilpotent leading coefficients can change
degree under multiplication, as the square-zero comparison illustrates.
\end{enumerate}

\clearpage
\section*{Verification and reproducibility}
\addcontentsline{toc}{section}{Verification and reproducibility}
The files in \texttt{verification/} use the Python standard library. The
primary graph check constructs each character's Laurent polynomial with
integer coefficients, reduces it modulo the appropriate cyclotomic
polynomial, and changes to the basis $1,(Z-1),\ldots,(Z-1)^{e-1}$. It
does not assume the case distinctions in the proofs. On small graphs
a second check constructs the reduced Laplacian directly and evaluates
its determinant by fraction-free elimination. This last check also tests
the parallel-edge convention at the bottom of the dyadic towers.
\begin{center}\small
\begin{tabular}{lp{.61\linewidth}}
\toprule
Program & Checks performed\\
\midrule
\texttt{check\_graphs.py} & All published DuBose--Vallières numerical levels, cofactors on graphs with at most 90 vertices in these cases, and six examples of the moment formula.\\
\texttt{check\_kataoka.py} & All nine Table 7 rows at $n=1,\ldots,9$; direct cofactors at levels 1--3.\\
\texttt{check\_dv3.py}, \texttt{check\_odd\_rows.py} & Higher-conductor tests through $n=12$ for the unequal-slope dyadic example and Kataoka's indices 5 and 7. The latter also checks the exceptional root-of-unity expansions.\\
\texttt{check\_products.py} & Counts disjoint polynomial chains, tests all 37 lines of a ternary cube of dimension three, and computes finite-factor density bounds exactly.\\
\bottomrule
\end{tabular}
\end{center}
The general statements about all levels follow from the proofs in the
text. The scripts provide both finite base cases and independent tests.
\texttt{source-audit.txt} gives pinned Lean paths, declaration lines, and
file hashes. \texttt{incoming-inventory.txt} records the 41 comparison
archives at the ProveIt revision inspected. The separate local-field note
in this package treats valuation losses over $\mathbb Q_3(\sqrt[3]{3})$.
\begin{thebibliography}{99}\small
\bibitem{Alpoge} L. Alpöge, \emph{Remarks on the Riemann hypothesis},
2026. \url{https://alpo.ge/remark.pdf}.
\bibitem{LeanRepo} L. Sernau and collaborators, \emph{GRH for Dirichlet
$L$-functions, in Lean 4}, 2026. Source, reading guide and verification
notes. \href{https://github.com/lsernau/grh-lean/tree/72fead102f663d5c7a407964a13dc0cf9d977207/grh}{Pinned readable edition}.
\bibitem{DV} S. DuBose and D. Vallières, \emph{On $\mathbb Z_\ell^d$-towers
of graphs}, Algebraic Combinatorics \textbf{6} (2023), 1331--1346.
Theorem 6.2 and Section 7. \url{https://doi.org/10.5802/alco.304}.
\bibitem{Kat} T. Kataoka, \emph{An Iwasawa-type asymptotic formula for
multiple $\mathbb Z_p$-coverings of graphs}, arXiv:2606.03579v1 (2026).
Section 6.3, Tables 6--7. \url{https://arxiv.org/abs/2606.03579}.
\bibitem{DLRV} C. Dion, A. Lei, A. Ray and D. Vallières, \emph{On the
distribution of Iwasawa invariants associated to multigraphs}, Nagoya
Math. J. \textbf{253} (2024), 48--90. \url{https://doi.org/10.1017/nmj.2023.18}.
\bibitem{KM} S. Kleine and K. Müller, \emph{On the growth of the Jacobians
in $\mathbb Z_p^l$-voltage covers of graphs}, Algebraic Combinatorics
\textbf{7} (2024), 1011--1038. \url{https://doi.org/10.5802/alco.366}.
\bibitem{Asada} M. Asada, E. Fourakis, S. Manski, N. McNew, S. J. Miller
and G. Moreland, \emph{Subsets of $\mathbb F_q[x]$ free of 3-term
geometric progressions}, Finite Fields Appl. \textbf{44} (2017), 135--147.
Question 6.5. \url{https://doi.org/10.1016/j.ffa.2016.10.002}.
\bibitem{Behrend} F. A. Behrend, \emph{On sets of integers which contain
no three terms in arithmetical progression}, Proc. Natl. Acad. Sci. USA
\textbf{32} (1946), 331--332. \url{https://doi.org/10.1073/pnas.32.12.331}.
\bibitem{DHJ} D. H. J. Polymath, \emph{A new proof of the density
Hales--Jewett theorem}, Ann. of Math. \textbf{175} (2012), 1283--1327.
\url{https://doi.org/10.4007/annals.2012.175.3.6}.
\bibitem{Moser} D. H. J. Polymath, \emph{Density Hales--Jewett and
Moser numbers}, arXiv:1002.0374. \url{https://arxiv.org/abs/1002.0374}.
\end{thebibliography}
